\chapter{Margin Models}\label{ch:rc:margin-models}

Over March 2020 the initial margin required by the world's clearing houses rose by roughly 300
billion dollars, and members posted another 115 billion of excess collateral on top: about 40\%
more than their February average, called in the weeks when every firm needed cash for something
else. The rise was not a failure of the models: it was what models that track volatility do when
volatility jumps. Bilateral margin under the industry's standard model barely moved, by design.
The difference between the two is a choice about procyclicality, and it decides who must find
liquidity in a crisis. One Quant Book~1 described margin and clearing; this chapter builds the
models: historical and filtered simulation at clearing houses, the sensitivity-based and
standardised approaches for bilateral trades, the tools that damp procyclicality, and the tests
that check a margin model covers what it should.

\section{What an initial-margin model must do}

\begin{definition}[Initial-margin model]\label{def:rc:margin-models:model}
An \emph{initial-margin model}\index{initial-margin model} sets the collateral a party must post
against the loss its counterparty would suffer in closing out its portfolio after a default, over
the margin period of risk, at a high confidence: at clearing houses at least 99\% single-tailed, for
bilateral trades 99\% over ten days.
\end{definition}

A margin model must be risk-sensitive (so that margin follows the portfolio), stable (so that
margin does not jump with every tick of volatility), transparent and replicable (so that members
can forecast their calls), and fast. The first two pull against each other.

\begin{example}[A member's swap book]\label{ex:rc:margin-models:book}
A member's cleared USD swaps have illustrative exposures of $+60$, $-40$, $-150$ and $-30$ thousand
dollars per basis point at two, five, ten and thirty years (the P\&L for a one-basis-point rise), a net
$-160$ thousand. With daily Treasury par yields standing in for swap rates, a five-day \omterm{def:rc:market-risk-measures:hs}{historical
simulation} at 99\% over one year of history gives an initial margin of USD~4.19 million on average in
February 2020.
\end{example}

\section{Clearing-house models: scenarios, historical and filtered simulation}

Clearing houses margin with \omterm{def:rc:market-risk-measures:hs}{historical simulation} over a look-back window, scenario grids (the
scanning ranges of One Quant Book~1), or \omterm{def:rc:market-risk-measures:fhs}{filtered historical simulation}, which rescales past moves
to current volatility (chapter~21).

\omcode{code/firm/initmargin/firm_initmargin.py}{23}{46}{Overlapping multi-day moves, historical margin, and filtering to today's volatility.}\label{lst:rc:margin-models:hs}

\begin{example}[March 2020 replayed]\label{ex:rc:margin-models:replay}
Recomputing the member's margin every day from January to June 2020, the one-year historical model
rises from USD~4.19 million (February average) to a peak of 5.25 million on 18 March, 25\% more. The
filtered model, with an EWMA decay of 0.97, rises from 4.21 million to 10.30 million on 23 March, 2.45
times its February level: USD~6.09 million of new margin in about three weeks
(\cref{fig:rc:margin-models:replay}). The historical model is slow because a one-year window dilutes a
few volatile days; the filtered model reacts at once.
\end{example}

\begin{omfigure}
\begin{tikzpicture}
\begin{axis}[width=11.5cm,height=5.6cm,
  xlabel={days from 2 January 2020},ylabel={initial margin (USD million)},
  tick label style={font=\small},label style={font=\small},
  xmin=0,xmax=180,ymin=0,ymax=12,grid=major,grid style={gray!25},
  legend columns=2,legend cell align=left,legend style={font=\footnotesize,at={(0.5,-0.25)},anchor=north,draw=none,fill=none,/tikz/every even column/.append style={column sep=8pt}},
  table/col sep=comma]
\addplot[omExo,very thick] table[x=day,y=hs]{figdata/rates-credit-risk/25-margin-models/im2020.csv};
\addplot[omThm,very thick] table[x=day,y=fhs]{figdata/rates-credit-risk/25-margin-models/im2020.csv};
\addplot[omDef,thick,dashed] table[x=day,y=buffer]{figdata/rates-credit-risk/25-margin-models/im2020.csv};
\addplot[omMeth,thick,dotted] table[x=day,y=stressed]{figdata/rates-credit-risk/25-margin-models/im2020.csv};
\legend{historical (1 year),filtered (EWMA 0.97),filtered with 25\% buffer,filtered with 25\% stressed weight}
\end{axis}
\end{tikzpicture}

\omcaption{Initial margin of the member's swap book through the first half of 2020 under four
models. Filtering to current volatility more than doubles the margin in March; a buffer or a
stressed-period weight raises the margin held before the stress and so trims the increase, without
changing the peak. Data: US Treasury par yields; the chapter's tutorial.}\label{fig:rc:margin-models:replay}
\end{omfigure}

\begin{dated}{2026-09}{Margin in March 2020}\label{dat:rc:margin-models:march2020}
The Basel Committee, CPMI and IOSCO reported in September 2022 that initial margin required across
clearing houses rose by roughly 300 billion dollars over March 2020, with 115 billion more of excess
collateral, about 40\% above the February average; bilateral margin under ISDA's standard model, used
by most banks under the margin rules, stayed relatively stable. About half of the clearing houses
surveyed had no formal anti-procyclicality framework; tolerances reported for margin increases ranged
from 25\% to 80\% over five days.
\end{dated}

\section{The standard bilateral model}

\begin{definition}[Standard initial margin model]\label{def:rc:margin-models:simm}
A \emph{standard initial margin model}\index{standard initial margin model} computes bilateral margin
from prescribed sensitivities (delta, vega, curvature) by risk class, weighted by calibrated risk
weights and aggregated with prescribed correlations, $\mathrm{IM} = \sqrt{WS^\top R\,WS}$ within each class: the
same structure as the FRTB sensitivities method (chapter~23). The industry's model is ISDA's SIMM,
whose parameters are recalibrated periodically; the chapter uses illustrative parameters.
\end{definition}

\begin{definition}[Schedule-based initial margin, initial-margin threshold]\label{def:rc:margin-models:schedule}
\emph{Schedule-based initial margin}\index{schedule-based initial margin} applies fixed shares of
notional by asset class and maturity (for interest rates 1\%, 2\% and 4\% for maturities up to two, two
to five and over five years; 6\% for FX; 15\% for equity), netted as $0.4\times\text{gross}+0.6\times\mathrm{NGR}\times\text{gross}$
with NGR the ratio of net to gross replacement cost. The \emph{initial-margin
threshold}\index{initial-margin threshold}, up to EUR~50 million per group, is the margin a party need
not collect.
\end{definition}

\begin{example}[Three bilateral numbers]\label{ex:rc:margin-models:bilateral}
Today, for the same book traded bilaterally: a ten-day historical 99\% margin of USD~4.15 million; a
sensitivity-based margin of 8.98 million with illustrative ten-day risk weights of 60, 60, 55 and 50
basis points; and a schedule margin of 13.19 million gross, 9.23 million net at an NGR of 0.5
(\cref{fig:rc:margin-models:approaches}). The schedule ignores the offsets between tenors; the
sensitivity model sees them only through its correlations.
\end{example}

\begin{omfigure}
\begin{tikzpicture}
\begin{axis}[width=11.5cm,height=5cm,xbar,bar width=12pt,
  symbolic y coords={schedule gross,schedule (NGR 0.5),sensitivity-based,HS 10-day},ytick=data,
  xlabel={initial margin (USD million)},
  tick label style={font=\small},label style={font=\small},
  xmin=0,xmax=15,grid=major,grid style={gray!25},
  nodes near coords,every node near coord/.append style={font=\footnotesize,/pgf/number format/fixed,/pgf/number format/precision=2},
  table/col sep=comma]
\addplot[fill=omDef!60,draw=omDef] table[y=approach,x=im]{figdata/rates-credit-risk/25-margin-models/approaches.csv};
\end{axis}
\end{tikzpicture}

\omcaption{Bilateral initial margin of the member's book today under a ten-day historical model,
the sensitivity-based model with illustrative parameters, and the regulatory schedule. Simpler
methods charge more. Data: US Treasury par yields; the chapter's tutorial.}\label{fig:rc:margin-models:approaches}
\end{omfigure}

\begin{dated}{2026-09}{Uncleared margin rules}\label{dat:rc:margin-models:umr}
Under the Basel Committee--IOSCO framework (as revised in April 2020), initial margin on
non-centrally cleared derivatives covers a one-tailed 99\% move over ten days; since 1 September 2022,
on a permanent basis, it applies between groups with more than EUR~8 billion of average aggregate
notional; the \omterm{def:rc:margin-models:schedule}{initial-margin threshold} is up to EUR~50 million and the \omterm{def:rc:counterparty-exposure:csa}{minimum transfer amount} up to
EUR~500\,000. ISDA's SIMM in force from 11 July 2026 is version 2.8 calibrated to December 2025, a
public document.
\end{dated}

\section{Procyclicality and its tools}

\begin{definition}[Anti-procyclicality tool]\label{def:rc:margin-models:apc}
An \emph{anti-procyclicality tool}\index{anti-procyclicality tool} limits the rise of margin in stress
by holding more in calm markets: a buffer above the model margin that can be released in stress, a
weight on a stressed period in the calibration, a long look-back window, or a floor on volatility.
European rules for central counterparties require at least one of a buffer of at least 25\% of the
calculated margin, a weight of at least 25\% on stressed observations, or margin no lower than a
ten-year look-back would give.
\end{definition}

\begin{definition}[Concentration add-on]\label{def:rc:margin-models:conc}
A \emph{concentration add-on}\index{concentration add-on} charges extra margin on positions large
relative to the market's capacity to absorb them, since closing them out would take longer than the
margin period of risk and move prices.
\end{definition}

\omcode{code/firm/initmargin/firm_initmargin.py}{73}{90}{Two anti-procyclicality tools: a buffer released in stress and rebuilt afterwards, and a stressed-period weight.}\label{lst:rc:margin-models:apc}

\begin{example}[The tools in March 2020]\label{ex:rc:margin-models:apc}
With a 25\% buffer on the filtered model, released when the book's EWMA volatility exceeds twice its
one-year average, the member held 5.26 million in February and 10.30 million at the peak: 5.04 million
of new margin instead of 6.09. With a 25\% weight on the margin of the most volatile year before 2020
(6.55 million), it held 4.79 million and had to find 5.51 million. The tools shift margin from the stress
to the calm; they do not lower the peak.
\end{example}

\section{Backtesting margin}

A margin model is backtested like a VaR model: count the days on which the realised loss over the
margin period of risk exceeded the margin held the day before. Clearing houses report their coverage
and investigate breaches.

\begin{example}[Coverage since 2017]\label{ex:rc:margin-models:backtest}
Over 2\,177 days since January 2017, the five-day loss of the member's book exceeded the one-year
historical margin 35 times and the filtered margin 37 times, against 22 expected at 99\%: both models
cover less than they aim to, largely in the 2020 and 2022 episodes, and overlapping windows make the
exceptions cluster.
\end{example}

\section{Tutorial: margin through a crisis}\label{tut:rc:margin-models:tutorial}

\begin{tutorial}
\textbf{Goal.} Compute the member's initial margin by historical and filtered simulation, replay the
first half of 2020 with and without \omterm{def:rc:margin-models:apc}{anti-procyclicality tools}, and compare the bilateral approaches.
\textbf{End state:} the numbers of
\cref{ex:rc:margin-models:replay,ex:rc:margin-models:bilateral,ex:rc:margin-models:apc,ex:rc:margin-models:backtest}
and the two charts.
\begin{tutsteps}
\item \textbf{Data}: \texttt{data()} reads the par yields; \texttt{pnl()} applies the exposures.
\item \textbf{Replay}: \texttt{replay\_2020()}.
\item \textbf{Bilateral}: \texttt{hs10\_today()}, \texttt{simm\_today()}, \texttt{schedule\_today()}.
\item \textbf{Backtest}: \texttt{backtest\_since()}; \texttt{fig\_rc\_margin.py} writes the charts.
\end{tutsteps}
\textbf{What to change next.} Use a ten-year look-back where history allows; add a volatility floor
at the ten-year average and compare the March increase.
\end{tutorial}

\section{Build: the initial-margin engine}\label{bld:rc:margin-models:initmargin}

\begin{build}
\bfield{Purpose} The firm's margin forecasts: what clearing houses and bilateral counterparties will
call, for liquidity planning, MVA (chapter~19) and \omterm{def:rc:stress-testing-and-scenarios:stress}{stress tests} (chapter~22).
\bfield{Interface} \texttt{overlapping}, \texttt{hs\_im}, \texttt{ewma\_vol}, \texttt{fhs\_moves};
\texttt{simm\_like(sens, tenors, rw)}; \texttt{schedule\_rate}, \texttt{schedule\_im(trades, net,
gross)}; \texttt{apc\_buffer}, \texttt{apc\_stressed\_weight}; \texttt{backtest}, \texttt{peak\_increase}.
\bfield{Rules} Losses positive; overlapping multi-day moves; illustrative sensitivity parameters;
the schedule and its NGR netting as in the Basel--IOSCO framework.
\bfield{Acceptance tests} \texttt{code/firm/initmargin/tests/}: overlapping sums; historical margin
of a known sample; schedule rates and netting; single-factor and offsetting sensitivity margins; buffer
and stressed-weight behaviour; backtest count.
\bfield{Stretch} Full-revaluation margin for options; concentration and liquidity add-ons; the licensed
SIMM with its calibration files; intraday margin.
\end{build}

\begin{omsources}
\item BCBS--CPMI--IOSCO, \emph{Review of margining practices}, September 2022.
\item BCBS--IOSCO, \emph{Margin requirements for non-centrally cleared derivatives}, April 2020.
\item CPMI--IOSCO, \emph{Principles for financial market infrastructures}, April 2012.
\end{omsources}

\section{Exercises}

\begin{exercise}[$\star$]\label{exo:rc:margin-models:1}
Compute the schedule margin of a USD~100 million ten-year swap and a USD~100 million one-year swap,
gross.
\end{exercise}

\begin{exercise}[$\star$]\label{exo:rc:margin-models:2}
Why is a one-year historical margin slow to react to a volatility spike?
\end{exercise}

\begin{exercise}[$\star$]\label{exo:rc:margin-models:3}
A counterparty's net replacement cost is 30 against a gross of 120. By what share does netting reduce
its schedule margin?
\end{exercise}

\begin{exercise}[$\star\star$]\label{exo:rc:margin-models:4}
Why did bilateral SIMM margin stay stable in March 2020 while CCP margin rose?
\end{exercise}

\begin{exercise}[$\star\star$]\label{exo:rc:margin-models:5}
Why do \omterm{def:rc:margin-models:apc}{anti-procyclicality tools} not lower the peak margin in the replay?
\end{exercise}

\begin{exercise}[$\star\star$]\label{exo:rc:margin-models:6}
What does a \omterm{def:rc:margin-models:conc}{concentration add-on} protect against that a volatility-based model misses?
\end{exercise}

\begin{exercise}[$\star\star\star$]\label{exo:rc:margin-models:7}
\emph{Coding.} Compute the expected number of exceptions and the Kupiec p-value for 35 exceptions in
2\,177 days at 1\%, treating days as independent, and say why the assumption is wrong.
\end{exercise}

\begin{exercise}[$\star\star\star$]\label{exo:rc:margin-models:8}
\emph{Find the flaw.} ``Our margin model is well calibrated: its February margin covered 99\% of
February's moves.''
\end{exercise}

\section{Problem: The Margin Calls of March 2020}

\begin{problem}[{Weekend problem --- finding the cash}]\label{pb:rc:margin-models:1}
The member of this chapter clears its swap book at a clearing house that uses the filtered model. Its
treasury holds liquidity for margin calls.

\medskip\noindent\textbf{Part I --- The calls.}
\begin{enumerate}
\item Give the February margin and the peak, with the date.
\item Give the cash to find over March.
\item What drove the increase: positions or volatility?
\item How much of the rise would the one-year historical model have called?
\item Why might the clearing house prefer the filtered model anyway?
\end{enumerate}

\medskip\noindent\textbf{Part II --- The tools.}
\begin{enumerate}[resume]
\item Give the February margin and the increase with a 25\% buffer.
\item Give them with a 25\% stressed-period weight.
\item What does each tool cost the member in calm years?
\item When should the buffer be rebuilt, and why slowly?
\item What reported tolerances for increases over five days did the 2022 review find?
\end{enumerate}

\medskip\noindent\textbf{Part III --- Bilateral.}
\begin{enumerate}[resume]
\item Give the member's bilateral margin under the three approaches.
\item Why would a member prefer to clear?
\item What is the effect of the EUR~50 million threshold on this book?
\item How does variation margin interact with initial margin in a crisis?
\item Which model would you use to forecast liquidity needs?
\end{enumerate}

\medskip\noindent\textbf{Part IV --- Judgement.}
\begin{enumerate}[resume]
\item Is procyclical margin a flaw or a feature?
\item Who bears the cost of anti-procyclicality?
\item How should a member size its liquidity buffer for margin?
\item State the \emph{named result}: the increase in initial margin with and without a 25\% buffer and a
stressed weight, and the cash the member had to find.
\item In one sentence: what does a margin model trade off?
\end{enumerate}
\end{problem}

\section{Interview questions}

\begin{interviewq}[$\star$ \iqroles{risk, trader}]\label{iq:rc:margin-models:1}
What is initial margin for, and how does a clearing house compute it?
\end{interviewq}

\begin{interviewq}[$\star\star$ \iqroles{risk}]\label{iq:rc:margin-models:2}
Explain procyclicality of margin and name three tools against it.
\end{interviewq}

\begin{interviewq}[$\star\star$ \iqroles{developer, risk}]\label{iq:rc:margin-models:3}
How does a sensitivity-based bilateral margin model work?
\end{interviewq}

\begin{interviewq}[$\star\star$ \iqroles{trader, bank}]\label{iq:rc:margin-models:4}
Why can two clearing houses charge different margin for the same swap?
\end{interviewq}

\begin{interviewq}[$\star\star\star$ \iqroles{researcher, risk}]\label{iq:rc:margin-models:5}
How would you backtest a margin model, and what makes it hard?
\end{interviewq}

\begin{interviewq}[$\star\star\star$ \iqroles{bank}]\label{iq:rc:margin-models:6}
What did March 2020 teach about margin and liquidity?
\end{interviewq}
